
%----------------------------------------------------------------------------------------
%	PACKAGES AND OTHER DOCUMENT CONFIGURATIONS
%----------------------------------------------------------------------------------------

\documentclass[12pt]{article} % Default font size is 12pt, it can be changed here
\usepackage{hyperref}
\usepackage{amsmath}
\usepackage{geometry} % Required to change the page size to A4
\geometry{a4paper} % Set the page size to be A4 as opposed to the default US Letter

\usepackage{graphicx} % Required for including pictures

\usepackage{float} % Allows putting an [H] in \begin{figure} to specify the exact location of the figure
\usepackage{wrapfig} % Allows in-line images such as the example fish picture

\usepackage{lipsum} % Used for inserting dummy 'Lorem ipsum' text into the template

\linespread{1.2} % Line spacing

%\setlength\parindent{0pt} % Uncomment to remove all indentation from paragraphs

\graphicspath{{./Pictures/}} % Specifies the directory where pictures are stored
\begin{document}

%----------------------------------------------------------------------------------------
%	TITLE PAGE
%----------------------------------------------------------------------------------------

\begin{titlepage}

\newcommand{\HRule}{\rule{\linewidth}{0.5mm}} % Defines a new command for the horizontal lines,
%change thickness here
\center % Center everything on the page
\includegraphics[width=\textwidth]{Glasgow}\\[1.5cm]
\textsc{\LARGE Algorithmics}\\[0.5cm] % Major heading such as course name
\textsc{\Large Assessed Exercise}\\[0.5cm] % Minor heading such as course title

\HRule \\[0.4cm]
{ \huge \bfseries Suffix Tree Applications}\\[0.4cm]
\\
\emph{State \& Strategy}
\HRule \\[1.5cm]

\begin{minipage}{0.4\textwidth}
\begin{flushleft} \large
\emph{Author:}\\
Garry \textsc{Sharp}\\
0801585s\\ % Your name
\end{flushleft}
\end{minipage}
~
\begin{minipage}{0.4\textwidth}
\begin{flushright} \large
\emph{Supervisors:} \\
Dr. D. \textsc{Manlove}\\ % Supervisor's Name
\end{flushright}
\end{minipage}\\[4cm]

{\large \today}\\[3cm] % Date, change the \today to a set date if you want to be precise

\vfill % Fill the rest of the page with whitespace

\end{titlepage}

\tableofcontents
\newpage
\section{Task 1 (Search Suffix Tree)}
The task here is implement a method \textit{searchSuffixTree} in the class \textit{SuffixTreeAppl}
which searches over a global instance of a \textit{SuffixTree} known as \textit{t} with the search
term being an array of bytes \textit{x}
\subsection{State}
My solution works in an effective manner to search the tree in \textit{O(n)} time. It should be
noted that the solutions to Task 2 \ref{task2}, Searching for all occurrences relies heavily on
this method. This means that there is a very small amount(almost negligable) amount of redundant
code when searching (sets a global variable). 
\subsection{Solution}
Firstly, lets examine the variables that are used. I keep track of an x index (initialised
to 0) a node index (index of the characted being compared at the node), a boolean called match
(which is set to true/false depending on whether there is a match), an integer start location
(initialised to -1) as well as a SuffixTreeNode called currentNode (initialised to the child of the
root of the tree).

My solution comparing the relevant values of x with the relevant values held at the node. If all
characters are matched successfully (and there are more in subsequent nodes) then I perform the
same check at the child of that node. If all characters have been matched and no more are left then
the loop terminates.

If there is a mismatch at any characters then I examine the sibling that node until, either there
is another match (is a mismatch occurs at that node it looks at the sibling, otherwise if fully
matched goes to the child) or a mismatch at which point the loop exits.

A variable t2Node is set throughout and this is so that Task 2 can be completed by recycling code,
this node is eventually set to the final node that contains leafs of the instances of the search
term.

If a match does occur and startLocation is not already set, a variable startLocation is set to it's
left label of the currentNode. Otherwise is a mismatch occurs startLabel is reset (set to -1).

The program terminates by setting the Task1Info object with position startLocation (which will be a
-1 if not found) to the main method which then displays it accordingly. The search executes in O(n)
time.

This solution uses no external helper methods.
\newpage
\section{Task 2 (All Occurrences)}
The task here is very similar to Task 1, but with the result being equal to every instance of the
search term in the text, specifically the starting index of every occurrence.
\label{task2}
\subsection{State}
This works in an effective manner by reusing the search functionality implemented in Task 1, this
is done by calling the searchSuffixTree method and then using the t2Node (which is set to the
branch node from which all decendent leaf nodes contain the occurrences). The solution executes in
O(n) time (plus build time).
\subsection{Solution}
As mentioned above, searchSuffixTree is called and t2Node is assigned accordingly. From here on I
used a recursive method on the child of the node until there are no more siblings. I basically look
at all nodes, where if there is no child (ie. node is a leaf) I add the value to a global variable
called occurrences, otherwise (if leaf) I recurse down the child of the currentnode and after move
to the sibling.

Finally I iterate through the value of occurences and add each of the suffixes of each node to the
Task2Info object which I then return.
\newpage
\section{Task 3 (Longest Repeating Substring)}
\subsection{State}
\subsection{Solution}
\newpage
\section{Task4 (Longest Common Substring)}
\subsection{State}
\subsection{Solution}
\newpage

%----------------------------------------------------------------------------------------
%	BIBLIOGRAPHY
%----------------------------------------------------------------------------------------

%\begin{thebibliography}{99} % Bibliography - this is intentionally simple in this template

%\bibitem[Figueredo and Wolf, 2009]{Figueredo:2009dg}
%Figueredo, A.~J. and Wolf, P. S.~A. (2009).
%\newblock Assortative pairing and life history strategy - a cross-cultural
%  study.
%\newblock {\em Human Nature}, 20:317--330.
 
%\end{thebibliography}

%----------------------------------------------------------------------------------------

\end{document}